% =====================================================================
%  Frontespizio ufficiale del Dipartimento di Fisica e Astronomia
%  Adattato dal modello frontespizio-fisica.tex.
%  Incluso da main.tex con % =====================================================================
%  Frontespizio ufficiale del Dipartimento di Fisica e Astronomia
%  Adattato dal modello frontespizio-fisica.tex.
%  Incluso da main.tex con % =====================================================================
%  Frontespizio ufficiale del Dipartimento di Fisica e Astronomia
%  Adattato dal modello frontespizio-fisica.tex.
%  Incluso da main.tex con % =====================================================================
%  Frontespizio ufficiale del Dipartimento di Fisica e Astronomia
%  Adattato dal modello frontespizio-fisica.tex.
%  Incluso da main.tex con \input{frontespizio}
%
%  Il modello impone \textwidth = 450pt e \oddsidemargin = 0pt: qui la
%  misura viene applicata con \newgeometry alla SOLA pagina del titolo e
%  revocata subito dopo con \restoregeometry, così che il resto della tesi
%  conservi i margini di 3 cm dichiarati nel preambolo.
%
%  I comandi \textcolor{red}{...} del modello servivano solo a evidenziare
%  i campi da compilare; una volta compilati equivalgono a testo nero e
%  sono stati rimossi, come prescrive il commento del modello stesso.
%
%  Rispetto al modello è stato eliminato il blocco del correlatore, che
%  non è previsto in questo lavoro.
% =====================================================================

\newgeometry{left=2.54cm, right=2.63cm, top=4.6cm, bottom=2cm}

\begin{titlepage}

\begin{center}
{{\Large{\textsc{Alma Mater Studiorum $\cdot$ Universit\`a di Bologna}}}}
\rule[0.1cm]{15.8cm}{0.1mm}
\rule[0.5cm]{15.8cm}{0.6mm}
\\\vspace{3mm}

{\small{\bf Dipartimento di Fisica e Astronomia ``Augusto Righi''\\
Corso di Laurea in Fisica}}

\end{center}

\vspace{23mm}

\begin{center}
%
% TITOLO DELLA TESI
%
{\LARGE{\bf Fra ordine e disordine}}\\[5mm]
{\Large{\bf Propriet\`a statistiche del linguaggio\\
generato da Large Language Model}}
\end{center}

\vspace{50mm} \par \noindent

\begin{minipage}[t]{0.47\textwidth}
%
% RELATORE
%
{\large{\bf Relatore: \vspace{2mm}\\
Prof. Mirko Degli Esposti}}
\end{minipage}
%
\hfill
%
\begin{minipage}[t]{0.47\textwidth}\raggedleft
{\large{\bf Presentata da:
\vspace{2mm}\\
%
% CANDIDATO
%
Samuele Zurlo}}
\end{minipage}

\vspace{40mm}

\begin{center}
%
% ANNO ACCADEMICO
%
Anno Accademico 2025/2026
\end{center}

\end{titlepage}

\restoregeometry

%
%  Il modello impone \textwidth = 450pt e \oddsidemargin = 0pt: qui la
%  misura viene applicata con \newgeometry alla SOLA pagina del titolo e
%  revocata subito dopo con \restoregeometry, così che il resto della tesi
%  conservi i margini di 3 cm dichiarati nel preambolo.
%
%  I comandi \textcolor{red}{...} del modello servivano solo a evidenziare
%  i campi da compilare; una volta compilati equivalgono a testo nero e
%  sono stati rimossi, come prescrive il commento del modello stesso.
%
%  Rispetto al modello è stato eliminato il blocco del correlatore, che
%  non è previsto in questo lavoro.
% =====================================================================

\newgeometry{left=2.54cm, right=2.63cm, top=4.6cm, bottom=2cm}

\begin{titlepage}

\begin{center}
{{\Large{\textsc{Alma Mater Studiorum $\cdot$ Universit\`a di Bologna}}}}
\rule[0.1cm]{15.8cm}{0.1mm}
\rule[0.5cm]{15.8cm}{0.6mm}
\\\vspace{3mm}

{\small{\bf Dipartimento di Fisica e Astronomia ``Augusto Righi''\\
Corso di Laurea in Fisica}}

\end{center}

\vspace{23mm}

\begin{center}
%
% TITOLO DELLA TESI
%
{\LARGE{\bf Fra ordine e disordine}}\\[5mm]
{\Large{\bf Propriet\`a statistiche del linguaggio\\
generato da Large Language Model}}
\end{center}

\vspace{50mm} \par \noindent

\begin{minipage}[t]{0.47\textwidth}
%
% RELATORE
%
{\large{\bf Relatore: \vspace{2mm}\\
Prof. Mirko Degli Esposti}}
\end{minipage}
%
\hfill
%
\begin{minipage}[t]{0.47\textwidth}\raggedleft
{\large{\bf Presentata da:
\vspace{2mm}\\
%
% CANDIDATO
%
Samuele Zurlo}}
\end{minipage}

\vspace{40mm}

\begin{center}
%
% ANNO ACCADEMICO
%
Anno Accademico 2025/2026
\end{center}

\end{titlepage}

\restoregeometry

%
%  Il modello impone \textwidth = 450pt e \oddsidemargin = 0pt: qui la
%  misura viene applicata con \newgeometry alla SOLA pagina del titolo e
%  revocata subito dopo con \restoregeometry, così che il resto della tesi
%  conservi i margini di 3 cm dichiarati nel preambolo.
%
%  I comandi \textcolor{red}{...} del modello servivano solo a evidenziare
%  i campi da compilare; una volta compilati equivalgono a testo nero e
%  sono stati rimossi, come prescrive il commento del modello stesso.
%
%  Rispetto al modello è stato eliminato il blocco del correlatore, che
%  non è previsto in questo lavoro.
% =====================================================================

\newgeometry{left=2.54cm, right=2.63cm, top=4.6cm, bottom=2cm}

\begin{titlepage}

\begin{center}
{{\Large{\textsc{Alma Mater Studiorum $\cdot$ Universit\`a di Bologna}}}}
\rule[0.1cm]{15.8cm}{0.1mm}
\rule[0.5cm]{15.8cm}{0.6mm}
\\\vspace{3mm}

{\small{\bf Dipartimento di Fisica e Astronomia ``Augusto Righi''\\
Corso di Laurea in Fisica}}

\end{center}

\vspace{23mm}

\begin{center}
%
% TITOLO DELLA TESI
%
{\LARGE{\bf Fra ordine e disordine}}\\[5mm]
{\Large{\bf Propriet\`a statistiche del linguaggio\\
generato da Large Language Model}}
\end{center}

\vspace{50mm} \par \noindent

\begin{minipage}[t]{0.47\textwidth}
%
% RELATORE
%
{\large{\bf Relatore: \vspace{2mm}\\
Prof. Mirko Degli Esposti}}
\end{minipage}
%
\hfill
%
\begin{minipage}[t]{0.47\textwidth}\raggedleft
{\large{\bf Presentata da:
\vspace{2mm}\\
%
% CANDIDATO
%
Samuele Zurlo}}
\end{minipage}

\vspace{40mm}

\begin{center}
%
% ANNO ACCADEMICO
%
Anno Accademico 2025/2026
\end{center}

\end{titlepage}

\restoregeometry

%
%  Il modello impone \textwidth = 450pt e \oddsidemargin = 0pt: qui la
%  misura viene applicata con \newgeometry alla SOLA pagina del titolo e
%  revocata subito dopo con \restoregeometry, così che il resto della tesi
%  conservi i margini di 3 cm dichiarati nel preambolo.
%
%  I comandi \textcolor{red}{...} del modello servivano solo a evidenziare
%  i campi da compilare; una volta compilati equivalgono a testo nero e
%  sono stati rimossi, come prescrive il commento del modello stesso.
%
%  Rispetto al modello è stato eliminato il blocco del correlatore, che
%  non è previsto in questo lavoro.
% =====================================================================

\newgeometry{left=2.54cm, right=2.63cm, top=4.6cm, bottom=2cm}

\begin{titlepage}

\begin{center}
{{\Large{\textsc{Alma Mater Studiorum $\cdot$ Universit\`a di Bologna}}}}
\rule[0.1cm]{15.8cm}{0.1mm}
\rule[0.5cm]{15.8cm}{0.6mm}
\\\vspace{3mm}

{\small{\bf Dipartimento di Fisica e Astronomia ``Augusto Righi''\\
Corso di Laurea in Fisica}}

\end{center}

\vspace{23mm}

\begin{center}
%
% TITOLO DELLA TESI
%
{\LARGE{\bf Fra ordine e disordine}}\\[5mm]
{\Large{\bf Propriet\`a statistiche del linguaggio\\
generato da Large Language Model}}
\end{center}

\vspace{50mm} \par \noindent

\begin{minipage}[t]{0.47\textwidth}
%
% RELATORE
%
{\large{\bf Relatore: \vspace{2mm}\\
Prof. Mirko Degli Esposti}}
\end{minipage}
%
\hfill
%
\begin{minipage}[t]{0.47\textwidth}\raggedleft
{\large{\bf Presentata da:
\vspace{2mm}\\
%
% CANDIDATO
%
Samuele Zurlo}}
\end{minipage}

\vspace{40mm}

\begin{center}
%
% ANNO ACCADEMICO
%
Anno Accademico 2025/2026
\end{center}

\end{titlepage}

\restoregeometry
